\pdfbookmark[1]{Abstract}{Abstract}
\chapter*{Abstract}

\textbf{The subject is what exact arithmetic lets a reading keep.} The anchor sift engine reaches a
symbol through equality alone, and a point either carries the same symbol as another or it does not,
with no order and no distance imposed from the reading end. On integers that equality is exact, and
an exact reading never merges two states that differ. This book reads the atom that way: an element
is the electrons of its neutral atom, each an integer quantum state, and the whole periodic table is
one accumulation of them.

\textbf{Pauli exclusion is an accumulation, and the reading keeps it.} Two fermions may not share a
state. Electrons accumulate into successive states instead of piling into the lowest, and that
stacking is the shell structure. The exclusion principle appears here as a count: the same
domain-blind primitive that finds two elements on one crystallographic site finds two electrons in
one state, and over all 118 elements it finds none. Coarsen the state and distinct electrons
agglomerate into one cell; the fullest cell at each coarsening is the Pauli capacity for that level,
read off the accumulation with nothing put in.

\textbf{The periodic law is the recurrence, measured against a drawn null.} An element's group
signature recurs down the table, and the recurrence is measured by an equality match rate read
against two backgrounds: a shuffle that keeps the counts and deletes the order, and a filling that
keeps no exclusion. The recurrence survives neither. It is in the order Pauli fills the shells and
not in the counts, and removing the exclusion leaves a ladder with no table. The row lengths, 8, 8,
18, 18, 32, 32, are read off the shell closures as the differences between them.

\textbf{What is measured and what is deferred.} The table pipeline runs on the ideal filling,
Madelung order under Hund's rule, and its figures are deterministic. A later chapter carries the Bohr
trajectory in exact rationals, where a radius is $n^2/Z$ and an energy is $-Z^2/n^2$ with nothing
rounded, and a whole spectral series reduces to line ratios that carry no measured constant at all.
The oracle chapter holds that ideal filling against the measured ground states from the NIST database,
covering all 118 elements, the measured configurations through element 108 and the predicted ones
above it, and reads off the aufbau exceptions where a real atom fills against the order. What is not
run is the comparison of an absolute Bohr line to a measured wavelength, which needs the recommended
constants and is named, not fetched. A final reading chapter takes the Standard Model particles, put
in from the particle data, for the structure their exact charges keep. The claims that rest on
measurement say so, and the claims that are exact arithmetic say that instead.
